\documentclass[
    headings=optiontotocandhead,% Erweiterung für das optionale Argument der
                                % Gliederungsbefehle aktiviert.
    twoside,
    numbers=noenddot,% Keine Punkte am Ende der Gliederungsnummern und davon
                     % abgeleiteten Nummern
    toc= %Flache TOC
    12pt, % Schriftgröße 
    titlepage, % es wird eine Titelseite verwendet 
    parskip=false, % Abstand zwischen Absätzen (ganze Zeile) 
    listof=totoc, % Verzeichnisse im Inhaltsverzeichnis aufführen 
    listof=flat, % mehr Abstand für grosse Zahlen
    numbers=noenddot, % kein Punkt am Ende bei N7ummern 
    %%enlargefirstpage,% Gibt es bei scrartcl nicht!!!!
    bibliography=totoc, % Literaturverzeichnis im Inhaltsverzeichnis aufführen 
    %index=totoc, % Index im Inhaltsverzeichnis aufführen 
    %captions=tableheading, % Beschriftung von Tabellen für Ausgabe oberhalb
                           % der Tabelle formatieren 
    %draft % Status des Dokuments (final/draft) draft hinzufügen zum anziegen 
    %%der zeilen ende
    a4paper,DIV=14,
    BCOR=15mm,
    % captions=tablesignature,
]{scrbook}

\setcounter{secnumdepth}{3}

\usepackage[T1]{fontenc}
\usepackage[utf8]{inputenc}

\usepackage[english, ngerman]{babel} % your native language must be the last one!!

\usepackage{lastpage}
\usepackage{listings}
\usepackage{blindtext}

%% Aufzählungen nicht so weit einrücken
\usepackage[inline]{enumitem}
%\setitemize{leftmargin=*} 
% Listen etwas wenige einrücken, erfordert enumitem
\setitemize{leftmargin=*}

\usepackage{lmodern}

\usepackage{xspace}

\usepackage{graphicx}

%%? \usepackage{textcomp}
\usepackage[hyphens]{url}
\usepackage{makeidx}
\makeindex
%%? \usepackage{graphicx}
\usepackage[numbers]{natbib}
\PassOptionsToPackage{normalem}{ulem}
\usepackage{ulem}

\usepackage{needspace}

\setlength\partopsep{0.5ex}%schoenere Listen
\usepackage[bottom]{footmisc}%fussnote ganz unten

\usepackage[]{microtype}
\UseMicrotypeSet[protrusion]{basicmath} % disable protrusion for tt fonts

\usepackage{multirow}   % Allows table elements to span several rows.
\usepackage{booktabs}   % Improves the typesettings of tables.
\usepackage{subcaption} % Allows the use of subfigures and enables their referencing.
\usepackage[ruled,linesnumbered,algochapter]{algorithm2e} % Enables the writing of pseudo code.
\usepackage[usenames,dvipsnames,table]{xcolor} % Allows the definition and use of colors. This package has to be included before tikz.
\usepackage{nag}       % Issues warnings when best practices in writing LaTeX documents are violated.
\usepackage{todonotes} % Provides tooltip-like todo notes.
\usepackage{float}
\usepackage{color}

%% bessere Suche im PDF
\input{glyphtounicode}
\pdfgentounicode=1
%%%%%%%%%%%%%%%%%%%%%%%%%%%%%%%%%%%%%%%%%%%%%%%%%%%%%%%%%%%%%%%%%%%%%%%%%%%%%%%%%%

%  Kopf und Fußzeilen -- links und rechts verschieden 
\newcommand{\kopfseitenummer}{{\bfseries \thepage}}
\newcommand{\kopfkapl}{{\bfseries\leftmark}}
\newcommand{\kopfkapr}{{\bfseries\rightmark}}
\newcommand{\kopfbild}{\voffset7mm\includegraphics[width=25mm]{HTL3RLogoRGB}}
\newcommand{\kopfHTL}{Höhere Technische Bundeslehranstalt Wien 3, \\Rennweg 	Abteilung für Informationstechnologie}

\usepackage[automark,headsepline,footsepline,plainfootsepline]{scrlayer-scrpage}
%\automark[chapter]{chapter}% Eventuell wenn doppelseitig
\setkomafont{pageheadfoot}{\normalcolor\footnotesize\scshape}
\setkomafont{pagenumber}{\normalfont\normalsize}
\clearpairofpagestyles
\ihead{\headmark}
\ohead{\kopfbild}
\ifoot{\kapitelautor}
\ofoot{\pagemark}
\ModifyLayer[addvoffset=-.6ex]{scrheadings.foot.above.line}% Linie verschieben
\ModifyLayer[addvoffset=-.6ex]{plain.scrheadings.foot.above.line}% Linie verschieben
\setlength{\headheight}{32pt}

% alle Seiten mit Kopfzeile
\renewcommand{\chapterpagestyle}{scrheadings}

%% Kapitel - aufwändige Kapitelüberschriften
%Options: Sonny, Lenny, Glenn, Conny, Rejne, Bjarne, Bjornstrup
%\usepackage[Bjornstrup]{fncychap}
% Alternative: 
%\usepackage{titlesec}

% Verzeichnisse - aufwändiger
%\usepackage{tocloft}


%% Code Beispiele
%% eine Variante 
\usepackage{listings}
\renewcommand{\lstlistingname}{\inputencoding{utf8}Listing}
%% andere Variante
%\usepackage{minted}
%\setminted{
%  linenos,
%  frame=lines,
%  framesep=2mm,
%  breaklines=true
%}
% Beispiel
%\begin{listing}[H]
%\begin{minted}{bash}
%...
%\end{minted}
%\caption{Beschreibung}
%\end{listing}
%% dritte Variante 
% mit/für pandoc
% create
% pandoc -s sample.md -o sample.tex
% -> copy/paste

\usepackage{fancyvrb}
\newcommand{\VerbBar}{|}
\newcommand{\VERB}{\Verb[commandchars=\\\{\}]}
\DefineVerbatimEnvironment{Highlighting}{Verbatim}{commandchars=\\\{\}}
% Add ',fontsize=\small' for more characters per line
\newenvironment{Shaded}{}{}
\newcommand{\KeywordTok}[1]{\textcolor[rgb]{0.00,0.44,0.13}{\textbf{#1}}}
\newcommand{\DataTypeTok}[1]{\textcolor[rgb]{0.56,0.13,0.00}{#1}}
\newcommand{\DecValTok}[1]{\textcolor[rgb]{0.25,0.63,0.44}{#1}}
\newcommand{\BaseNTok}[1]{\textcolor[rgb]{0.25,0.63,0.44}{#1}}
\newcommand{\FloatTok}[1]{\textcolor[rgb]{0.25,0.63,0.44}{#1}}
\newcommand{\ConstantTok}[1]{\textcolor[rgb]{0.53,0.00,0.00}{#1}}
\newcommand{\CharTok}[1]{\textcolor[rgb]{0.25,0.44,0.63}{#1}}
\newcommand{\SpecialCharTok}[1]{\textcolor[rgb]{0.25,0.44,0.63}{#1}}
\newcommand{\StringTok}[1]{\textcolor[rgb]{0.25,0.44,0.63}{#1}}
\newcommand{\VerbatimStringTok}[1]{\textcolor[rgb]{0.25,0.44,0.63}{#1}}
\newcommand{\SpecialStringTok}[1]{\textcolor[rgb]{0.73,0.40,0.53}{#1}}
\newcommand{\ImportTok}[1]{#1}
\newcommand{\CommentTok}[1]{\textcolor[rgb]{0.38,0.63,0.69}{\textit{#1}}}
\newcommand{\DocumentationTok}[1]{\textcolor[rgb]{0.73,0.13,0.13}{\textit{#1}}}
\newcommand{\AnnotationTok}[1]{\textcolor[rgb]{0.38,0.63,0.69}{\textbf{\textit{#1}}}}
\newcommand{\CommentVarTok}[1]{\textcolor[rgb]{0.38,0.63,0.69}{\textbf{\textit{#1}}}}
\newcommand{\OtherTok}[1]{\textcolor[rgb]{0.00,0.44,0.13}{#1}}
\newcommand{\FunctionTok}[1]{\textcolor[rgb]{0.02,0.16,0.49}{#1}}
\newcommand{\VariableTok}[1]{\textcolor[rgb]{0.10,0.09,0.49}{#1}}
\newcommand{\ControlFlowTok}[1]{\textcolor[rgb]{0.00,0.44,0.13}{\textbf{#1}}}
\newcommand{\OperatorTok}[1]{\textcolor[rgb]{0.40,0.40,0.40}{#1}}
\newcommand{\BuiltInTok}[1]{#1}
\newcommand{\ExtensionTok}[1]{#1}
\newcommand{\PreprocessorTok}[1]{\textcolor[rgb]{0.74,0.48,0.00}{#1}}
\newcommand{\AttributeTok}[1]{\textcolor[rgb]{0.49,0.56,0.16}{#1}}
\newcommand{\RegionMarkerTok}[1]{#1}
\newcommand{\InformationTok}[1]{\textcolor[rgb]{0.38,0.63,0.69}{\textbf{\textit{#1}}}}
\newcommand{\WarningTok}[1]{\textcolor[rgb]{0.38,0.63,0.69}{\textbf{\textit{#1}}}}
\newcommand{\AlertTok}[1]{\textcolor[rgb]{1.00,0.00,0.00}{\textbf{#1}}}
\newcommand{\ErrorTok}[1]{\textcolor[rgb]{1.00,0.00,0.00}{\textbf{#1}}}
\newcommand{\NormalTok}[1]{#1}


%% should be last packages
\usepackage{scrhack}

\usepackage[unicode=true,
 bookmarks=true,bookmarksnumbered=false,bookmarksopen=false,
 breaklinks=true,pdfborder={0 0 0},backref=false,colorlinks=false]
 {hyperref}
\hypersetup{pdftitle={Insight},
 pdfauthor={Wer auch immer},
 pdfsubject={Diplomarbeit},
 pdfkeywords={dies, das}}
\urlstyle{same} % don't use monospace font for urls

%% for pandoc
\providecommand{\tightlist}{%
  \setlength{\itemsep}{0pt}\setlength{\parskip}{0pt}}

% Auch Fußnoten bündig ausrichten
\deffootnote[]{1em}{1em}{\textsuperscript{\thefootnotemark\ }}
%% setup
\sloppy % weniger Meldungen
\voffset7mm % etwas nach unten

%% schöner: 10000 -- gar keine, 1000 als Mittelweg
\clubpenalty = 1000 % Schusterjungen verhindern
\widowpenalty = 1000 % Hurenkinder verhindern
\displaywidowpenalty = 1000 

%%%%%%%%%%%%%%%%%%%%%%%%%%%%%%%%%%%%%%%%%%%%%%%%%%%%%%%%%%%%%%%%%%%%%%%%%%%%%%%%%%
\begin{document}
%% wir schreiben keine Umlaut mit "a "o
\shorthandoff{"}
%% mit kapitelautor kann man den Autor festlegen oder auf leer setzen - steht dann in der Fußzeile.
\newcommand{\kapitelautor}{}

%%
\newcommand{\strong}[1]{\textbf{#1}}
\newcommand{\code}[1]{\texttt{#1}}

% einfaches "siehe ..." - das Ziel muss man markieren mit \label{name} -- macht pandoc automatisch
\newcommand{\kap}[1]{Kapitel~\ref{#1}, Seite~\pageref{#1}}
\newcommand{\siehe}[1]{siehe \kap{#1}}
\newcommand{\abb}[1]{Abbildung~\ref{#1}, Seite~\pageref{#1}}

%% http://ieg.ifs.tuwien.ac.at/~aigner/download/tuwien.sty
%Div. Abkürzungen (in Anlehnung an Jochen Köpper, jkthesis):
%\RequirePackage{xspace}
\newcommand{\bzw}{bzw.\@\xspace}
\newcommand{\bzgl}{bzgl.\@\xspace}
\newcommand{\ca}{ca.\@\xspace}
\newcommand{\dah}{d.\thinspace{}h.\@\xspace}
\newcommand{\Dah}{D.\thinspace{}h.\@\xspace}
\newcommand{\ds}{d.\thinspace{}s.\@\xspace}
\newcommand{\evtl}{evtl.\@\xspace}
\newcommand{\ua}{u.\thinspace{}a.\@\xspace}
\newcommand{\Ua}{U.\thinspace{}a.\@\xspace}
\newcommand{\usw}{usw.\@\xspace}
\newcommand{\va}{v.\thinspace{}a.\@\xspace}
\newcommand{\vgl}{vgl.\@\xspace}
\newcommand{\zB}{z.\thinspace{}B.\@\xspace}
\newcommand{\ZB}{Zum Beispiel\xspace} 

%% https://github.com/Digital-Media/HagenbergThesis
\newcommand{\latex}{La\-TeX\xspace} % kein schnoerkeliges LaTeX mehr
\newcommand{\tex}{TeX\xspace}       % kein schnoerkeliges TeX mehr
\newcommand{\bs}{\textbackslash}    % Backslash character
\newcommand{\obnh}{\hskip 0pt } %optional break without hyphen: e.g. PlugIn{\obnh}Filter

\newcommand{\sa}{s.\ auch\@\xspace}
\newcommand{\so}{s.\ oben\xspace}
\newcommand{\su}{s.\ unten\@\xspace}

\newcommand{\uae}{u.\thinspace{}\"A.\@\xspace}
\newcommand{\uva}{u.\thinspace{}v.\thinspace{}a.\@\xspace}
\newcommand{\uvm}{u.\thinspace{}v.\thinspace{}m.\@\xspace}



%%%%%Anfang Titelseite
%\pagenumbering{roman}
\frontmatter % Switches to roman numbering
\title{Diplomarbeit}
\begin{titlepage}
\begin{minipage}[b]{1\columnwidth}
\parbox[b]{50mm}{\includegraphics[width=45mm]{HTL3RLogoRGB}}
\hfill
\parbox[b]{130mm}{\footnotesize \textsc{Höhere Technische Bundeslehranstalt} Wien 3, Rennweg\\
IT \& Mechatronik\\
\\
HTL Rennweg :: Rennweg 89b\\
A-1030 Wien :: Tel +43 1 24215-10 :: Fax DW 18
}\\
\mbox{}
\end{minipage}

\vspace{1cm}


\begin{center}
\textbf{\LARGE{}Diplomarbeit}{\large{}}\\
\textbf{\large{}Insight}\\
 \vspace{15mm}
 ausgeführt an der\\
 Höheren Abteilung für Informationstechnologie/Medientechnik\\
 der Höheren Technischen Lehranstalt Wien 3 Rennweg\\
 \vspace{1cm}
 im Schuljahr 2017/2018\\
 \vspace{1cm}
 durch\\
 \vspace{0.5cm}
\textbf{\large{}Akyokus Hatice}\\
\textbf{\large{}Kienreich Niklas}\\
\textbf{\large{}Schön Kerstin}\\

\par\end{center}{\large \par}

\begin{center}
\vspace{20mm}
 \normalsize unter der Anleitung von\\
 \vspace{0.5cm}
 Mag. Roman Jerabek\\
 Mag. Andreas Fink
\par\end{center}

\begin{center}
\vspace{5mm}
Wien, \today 
\par\end{center}

\end{titlepage}%%%%%%%%%%%%%%%%%%%%% Ende Titelseite %%%%%%%%%%%%%%%%%%%%%%

\chapter*{Kurzfassung}

Was ist die Medientechnik? Was wird davon an der HTL beigebracht? \\
Mit diesen Fragen beschäftigen sich Schüler, die sich für diese Schule interessieren. Dabei haben Sie nur eine Möglichkeit sich darüber zu informieren, den Tag der offenen Tür. Am Tag der offenen Tür stehen die drei Bereiche Foto, Video und Audio im Vordergrund, wobei Webdesign, Softwareentwicklung und Webtechnik, nur kurz oder gar nicht erwähnt werden. Das große Problem hierbei ist, dass sich Interessenten an der Schule anmelden und merken, dass der Schwerpunkt bei der Programmierung liegt. Um dieses Problem zu vermeiden, entwickeln wir eine Webplattform, welche es Schülern ermöglicht sich über die Medientechnik zu informieren. Auf der Webseite befinden sich Videos sowie interaktive Elemente. Mittels Interviews werden wichtige Fragen über die Abteilung, sowie über die Schule im Allgemeinen beantwortet. Der User soll aber auch die Möglichkeit bekommen, medientechnische Elemente zu betrachten und auszuprobieren. Hierbei spielen interaktive Elemente eine große Rolle, da sie einerseits die User neugierig machen aber auch Beispiele sind, die tatsächlich beigebracht werden. Somit können sich Interessenten ein Bild von der Medientechnik machen, ohne das Missverständnisse entstehen.

\chapter*{Abstract}
\selectlanguage{english}

Thats why -- the translated text ,,Kurzfassung`` (this should be
a translation).

Im englischen Abstract sollte inhaltlich das Gleiche stehen wie in
der deutschen Kurzfassung. Versuchen Sie daher, die Kurzfassung präzise
umzusetzen, ohne aber dabei Wort für Wort zu übersetzen. Beachten
Sie bei der Übersetzung, dass gewisse Redewendungen aus dem Deutschen
im Englischen kein Pendant haben oder völlig anders formuliert werden
müssen und dass die Satzstellung im Englischen sich (bekanntlich)
vom Deutschen stark unterscheidet. Es empfiehlt sich übrigens – auch
bei höchstem Vertrauen in die persönlichen Englischkenntnisse – eine
kundige Person für das „proof reading“ zu engagieren. Die richtige
Übersetzung für „Diplomarbeit“ ist übrigens schlicht thesis, allenfalls
„diploma thesis“ oder „Master’s thesis“, auf keinen Fall aber „diploma
work“ oder gar „dissertation“\citep{hagenberg}.

Wichtig ist wegen des Abteilens ein \code{\textbackslash{}begin\{english\}}
bzw. \code{\textbackslash{}selectlanguage\{ngerman\}}.

\selectlanguage{ngerman}

\chapter*{Ehrenwörtliche Erklärung}

Ich erkläre an Eides statt, dass ich die individuelle Themenstellung
selbstständig und ohne fremde Hilfe verfasst, keine anderen als die
angegebenen Quellen und Hilfsmittel benutzt und die den benutzten
Quellen wörtlich und inhaltlich entnommenen Stellen als solche erkenntlich
gemacht habe.

\begin{flushleft}
\bigskip{}
Wien, am \today \\
\newcommand{\namesigdate}[2][8cm]{
\vspace{2cm}~\newline
\parbox{#1}{\hrule\centering #2\Large\strut}
\hfill
}
\namesigdate{Mitarbeiter Eins}
\namesigdate{Mitarbeiter Zwei} 
\namesigdate{Mitarbeiter Drei}
\par\end{flushleft}



%%%%%%%%%%%%%%%%%%%%%%%%%%%%%%%%%%%%%%%%%%%%%%%%%%%%%%%%%%%%%%%%%%%%%%%%%%%%%%%%%%%%%%%%
\cleardoublepage{}
\setcounter{tocdepth}{2}
\tableofcontents{}
\cleardoublepage{}
\listoftables
\cleardoublepage{}
\listoffigures

%hier geht es los mit dem Text - auf einer rechten Seite
\cleardoublepage{}
%\pagenumbering{arabic}
\mainmatter







% Das komplette nächste Kapitel wird in der externen Datei diplomarbeit2.tex gespeichert. 
% Es wird an dieser Stelle im Dokument eingebaut.
% Damit ist es möglich, mehrere Personen an diverse Teile der Diplomarbeit arbeiten zu lassen.
\chapter{Das Team}
\section{Hatice Akyokus}
\section{Niklas Kienreich}
\section{Kerstin Schön}
\chapter{Das Projekt}
\renewcommand{\kapitelautor}{Autor: Hatice Akyokus}
In diesem Kapitel wird die Planung des Projektes erläutert. 

\section{Die Idee}
Am Tag der offenen Tür werden jungen Schülern gezeigt, was die Medientechnik beinhaltet. Dabei werden nur die Bereiche Video, Audio und Fotografie demonstriert. Bereiche wie Webdesign, Softwaretechnik und Webtechnologie werden gänzlich ausgelassen. Dadurch, dass die Besucher denken, dass es nur diese multimedialen Bereiche gibt, melden sie sich an und merken schnell, dass man sich erst ab der 4. Klasse auf diese Bereiche spezialisiert. Um dies zu verhindern, wurden Informationsvideos gedreht, die auf einer Webseite zur Verfügung stehen. Es wurden insgesamt fünf Videos auf die Webseite gestellt. Zwei davon sind Interviews mit dem Abteilungsvorstand der HTL Rennweg und dem Medientechniklehrer Mag. Roman Jerabek. Die übrigen Videos sorgen für den „Fun-Factor“ und sind humorvoller als die Interviews. Die Webseite beinhaltet interaktive Elemente, welche zeigen, was mit der Medientechnik gemacht werden kann. Durch diese Mittel soll gezeigt werden, was in den fünf Jahren beigebracht wird.


\section{Die Planung}

Die Projektplanung ist die wichtigste Phase eines Projektes und erfolgt nach der Definierung der Projektziele.\footnote{\label{foot:2}vgl.https://www.rechnungswesen-verstehen.de/lexikon/projektplanung.php [Zugriff: 17.03.2018]}   Weiters wird die Projektmanagementmethode und alle, für diese Methode wichtigen, Tools, welche verwendet wurden, erklärt. Im Rahmen dieser Diplomarbeit, wurde als Erstes ein Kick-Off-Meeting mit den Projektbetreuern organisiert, bei der die Durchführung besprochen wurde. Anschließend wurden, die zu erledigenden Dokumente, die Umwelt- bzw. Umfeldanalyse und die Risikoanalyse, festgelegt. Zuletzt wurde die Projektmanagementmethode festgelegt, welche für die Diplomarbeit am besten geeignet ist. Nach dem Kick-Off-Meeting wurde der Antrag, welcher die Projektziele, die Projektorganisation, die Umwelt- bzw. Umfeldanalyse, die Risikoanalyse, die Meilensteinliste, die Projektressourcen und externe Kooperationspartner enthält, geschrieben und dem Abteilungsvorstand der Informationstechnologie, Dr. Gerhard Hager, übergeben.

\section{Projektmanagement und Kommunikation}

Für diese Diplomarbeit entschied sich das Team dazu, die Projektmanagementmethode Scrum zu verwenden. Für diese Methode ist ein Tool, welches die Arbeitspakete enthält, notwendig. Hierbei verwendet das Team die Webseite taiga.io, welches mit der Webseite für Zeiterfassung, toggl.com, verwendet werden kann. Für die Dokumentenverwaltung wird Sharepoint verwendet, welche mit den HTL Rennweg Accounts verbunden sind. Dabei hat der Hauptbetreuer, sowie der Nebenbetreuer Zugriff darauf und kann alle Dokumente einsehen.

\subsection[Projektmanagementsmethode]{Projektmanagementmethode\protect\footnote{\label{foot:2}vgl.http://projektmanagement-definitionen.de/glossar/scrum/  [Zugriff: 17.03.2018]}}

Die vom Projektteam gewählte Projektmanagementmethode, ist Scrum. Scrum ist eine agile Projektmanagementmethode und wird seit den 90er Jahren in der Softwareentwicklung verwendet und findet dort seinen Ursprung. Es werden drei Rollen definiert, der Product Owner, das Entwicklungsteam und der Scrum Master. 

\subsubsection[Rollen]{Rollen\protect\footnote{\label{foot:2}vgl.https://www.scrum.de/was-macht-product-owner/ [Zugriff: 17.03.2018]} \protect\footnote{\label{foot:2}vgl.https://www.projektmagazin.de/glossarterm/entwicklungsteam-scrum [Zugriff: 17.03.2018]} \protect\footnote{\label{foot:2}vgl.http://agiles-projektmanagement.org/scrum-rollen/scrum-master/ [Zugriff: 17.03.2018]}}   
\paragraph{Product Owner}
\leavevmode \\
Der Product Owner stellt Anforderungen an das Produkt und priorisiert diese. Dabei handelt es sich immer um eine Einzelperson und keinen Ausschuss. Somit darf diese Rolle auch nicht aufgeteilt werden. Weiters, ist er für das Product Backlog zuständig und sorgt dafür, dass der Wert des Produkts maximiert wird, in dem er die Einträge maximiert und so definiert, dass das Entwicklungsteam die Einträge versteht. Er ist die einzige Person, welche Änderungen vornehmen darf und dafür rechenschaftspflichtig ist. Die Entscheidung, ob ein Sprint vorzeitig beendet werden darf, liegt ebenfalls beim Product Owner. Doch die Hauptaufgabe des Product Owners liegt darin, Rücksprache mit dem Management und Stakeholdern zu halten, um die Arbeit zu optimieren und somit eine Fortschrittskontrolle zur Erreichung des Sprintziels durchzuführen. Im Sprintreview erklärt er, welche Product-Backlog-Einträge fertig sind und welche noch zu erledigen sind. 

\paragraph{Entwicklungsteam}
\leavevmode \\
Alle Mitglieder des Entwicklungsteams sind gleichberechtigt und haben untereinander keine Hierarchie. Somit haben sie keine Weisungsbefugnisse und tragen zur Erstellung des Produktes bei. Mitglieder ohne IT-Kenntnisse werden ebenfalls als Entwickler bezeichnet. Allerdings sollte das Entwicklungsteam eine Größe zwischen drei und neun Personen betragen. Die vom Product Owner spezifizierten Einträge, werden vom Entwicklungsteam umgesetzt, wobei sie die Machbarkeit und den Aufwand der Backlog Einträge beim Sprint Planning beurteilen und über den Umfang des Backlogs entscheiden. Das Entwicklungsteam entwickelt das Produkt und ist für die Abarbeitung des Backlogs zuständig. Allerdings in einer Reihenfolge, welche vom Product Owner festgelegt wird. Am Ende eines jeden Sprints wird ein Prototyp präsentiert und über den erzielten Fortschritt berichtet. 

\paragraph{Scrum Master}
\leavevmode \\
Der Scrum Master ist kein Projektleiter, sondern dafür verantwortlich, dass die Projektmanagementmethode Srum, richtig ausgeführt wird. Falls Scrum neu eingeführt wird, ist es seine Aufgabe und Pflicht, dem restlichen Team den Ablauf zu erklären. Er trifft während des Projektablaufes keine Entscheidungen und darf nicht in Entwicklungsprozess eingreifen. Sein Aufgabenbereich liegt in der Moderation und der Vermittlung an Informationen zwischen Product Owner und dem Entwicklungsteam. Ein Scrum Master entlastet weiters das Projektteam, in dem er zentrale Aufgaben übernimmt und über die Werte und Regeln eines Projektes wacht. Weiters leitet er Meetings, in welchen er unparteiisch zu handeln hat. Somit hat ein Scrum Master Meinungsverschiedenheiten oder generelle Konflikte zu beseitigen. 

\subsubsection[Product Backlog]{Product Backlog\protect\footnote{\label{foot:2}vgl.http://scrum-fibel.de/artefakte/sprint\%20backlog.html [Zugriff: 17.03.2018}}
Der Product Backlog wird vom Product Owner angelgt, gepflegt, geordnet und priorisiert und ist eine Sammlung von Anforderungen. Das Backlog enthält „User-Stories“, welche die Bedürfnisse der Nutzer widerspiegeln. Diese enthalten grobe Anforderungen, die für das Endergebnis wichtig sind und werden beim Sprint-Planning Meeting festgelegt. Die User-Stories müssen in kleinstmögliche Aufgaben aufgegliedert werden, damit das Team diese im Sprint abarbeiten kann. Ein Sprint dauert in der Regel ein bis vier Wochen, diese wird vom Scrum Master festgelegt. Weiters spielt die Priorisierung eine wichtige Rolle, denn diese entscheidet darüber, was im nächsten Sprint abgearbeitet wird. Diese darf während dem Sprint nicht verändert werden. Allerdings werden sie in Absprache mit den Stakeholdern korrigiert und angepasst, sofern die Dringlichkeit und Notwendigkeit eine große Rolle spielen. Um einen Überblick über alle Aufgaben zu behalten, wird ein Sprint Backlog angelegt. 

\subsubsection[Sprint Backlog]{Sprint Backlog\protect\footnote{\label{foot:2}vgl.http://scrum-fibel.de/artefakte/sprint\%20backlog.html [Zugriff: 17.03.2018]}}
In einem Sprint Backlog befindet sich eine Liste von Aufgaben. Diese sind essentiell um die Anforderungen des Product Backlogs umzusetzen. Das Sprint Backlog entsteht aus der Sprint Planung, wobei der Aufwand einer Aufgabe in Stunden geschätzt wird. Die Aufgaben nimmt sich jedes Teammitglied selbst und nutzt den Sprint Backlog zur Verfolgung der Aufgaben. 

\subsubsection[Sprint Planning]{Sprint Planning\protect\footnote{\label{foot:2}vgl.http://doktor-scrum.de/sprint-planning-1/ [Zugriff: 17.03.2018]}}
Im Sprint Planning werden die Anforderungen und Arbeitspakete, die im aktuellen Sprint umgesetzt werden sollen, besprochen. Der Product-Owner organisiert alles, das Entwicklungsteam, der Scrum-Master und optional der Stakeholder, nehmen teil.  Das Sprint Planning findet zu Beginn jedes Sprints statt und dauert vier bis acht Stunden. Aus dem Product-Backlog werden eine Anzahl an höchstpriorisierten Anforderungen ausgewählt und in ein Selected-Backlog übernommen. Diese sind nur Vorschläge, denn das Entwicklerteam entscheidet letztendlich, was im Sprint realisiert werden soll. Die Sprint-Planung teilt sich in zwei Phasen auf. In der ersten Phase, werden die Anforderungen detailliert besprochen. Dabei schätzt das Entwicklungsteam den Realisationsaufwand und legt fest, was im Sprint umgesetzt wird. In der zweiten Phase bespricht das Entwicklungsteam die Realisierung und zerlegt die Anforderungen in einzelne Tasks. Bei Fragen sollte der Product Owner erreichbar sein.


\subsubsection[Daily Scrum]{Daily Scrum\protect\footnote{\label{foot:2}vgl.http://scrum-master.de/Scrum-Meetings/Daily\underline Scrum\underline Meeting [Zugriff: 17.03.2018]}}
Ein Daily Scrum-Meeting ist ein kurzes, tägliches Meeting des Teams, welches ungefähr 15 Minuten dauert. Der Scrum Master nimmt teil und greift moderierend ein, wenn nötig. Beim Daily Scrum werden Statusupdates und Geschehnisse geschildert. 

\subsubsection[Sprint Review]{Sprint Review\protect\footnote{\label{foot:2}vgl.https://www.projektmagazin.de/methoden/sprint-review [Zugriff: 18.03.2018]}}
Das Sprint Review-Meeting findet am Ende jedes Sprintes statt. Sie dient zur Einholung von Feedback von den Stakeholdern. Dabei wird das Produktinkrement vorgestellt, wobei der Stakeholder seine Meinung gibt und dadurch das Produkt im nächsten Sprint optimiert werden kann. 

\subsection{Tool für Scrum}
Taiga ist ein Open Source Tool, welches für Scrum Projekte verwendet werden kann. Durch das schlichte und moderne Design, findet man sich schnell zurecht. Man bekommt die Möglichkeit geboten entweder ein Scrum oder Kanban Template zu benutzen. In dieser Diplomarbeit wurde das Scrum Template verwendet. Taiga bietet ein Interface für das Product und Sprint Backlog und dem dazugehörigen Taskboard, welche alle Tasks grafisch darstellt und sie mit einem Status versieht, um ersichtlich zu machen, in welcher Phase sich der Task befindet. Es gibt fünf Phasen: 
\begin{itemize}
\item Neu
\item In Arbeit
\item Bereit zum Testen
\item Fertig
\item Informationen werden benötigt
\end{itemize}

Die Tasks können in die jeweilige Phase verschoben werden und sobald sie sich in der Phase „Fertig“ befinden, aktualisiert sich der Projektfortschritt. Das gibt einen genauen Überblick über den Stand des Projektes. 

\subsection{Zeiterfassung}
Für die Zeiterfassung wurde toggl.com verwendet. Durch die Eingabe des Tasks und den Klick eines Buttons, wird die Zeit aufgezeichnet und in den Reports aufgelistet. Der Projektleiter kann die erfassten Zeiten jedes Teammitgliedes sehen und auch verändern. Es gibt zwei Möglichkeiten die Zeit zu erfassen:
\begin{itemize}
\item Manual Mode
\item Timer Mode
\end{itemize}

Im Manual Mode wird die Zeit händisch eingetragen, währenddessen im Timer Mode die Zeit, ähnlich wie bei einer Stoppuhr, automatisch läuft. Das führt zu einer genaueren Zeiterfassung. 

\subsection{Kommunikationsmittel}
In der Diplomarbeit wurde teamintern WhatsApp als Kommunikationsmittel verwendet. WhatsApp ist eine Instant-Messaging-App mit der Nachrichten in Echtzeit ausgetauscht werden können. Die Projektleitung erstellte zu diesem Zwecke eine Gruppe, in welcher alle Mitglieder eingeladen wurden. In dieser Gruppe wurden wichtige Ereignisse besprochen, sowie Fragen geklärt. 
Mit den Diplomarbeitsbetreuern wurde die Kommunikation durch E-Mail bevorzugt. 

\subsection{Dokumentenverwaltung durch Sharepoint}
Für die Dokumentenverwaltung, entschied sich das Team für Microsoft Sharepoint. Dieses Tool bietet eine Verbesserung der Teamproduktivität, sowie ein bequemes Verwalten von Dokumente durch eine kurze Einarbeitungszeit und einer schnellen und leichten Erstellung einer Zusammenarbeitsumgebung. Auf Sharepoint haben das gesamte Team und die Projektbetreuer Zugriff, um hochgeladene Dateien bewerten zu können.


\chapter{Ziele}
\renewcommand{\kapitelautor}{Autor: Hatice Akyokus}
Ziel unserer Diplomarbeit ist es, Jugendlichen der Unterstufe die Medientechnik näherzubringen. Dabei greift das Team auf klassische Storytelling-Methoden zurück, welche die Spannung, und den Spaß erhöhen sollen.\\ Folgende Ziele wurden im Diplomarbeitsantrag definiert:

\section{Muss-Ziele}
\begin{enumerate}
\item RE-M 1 Webseite\\
Eine Webseite ist erstellt.\\
Auf der Webseite kann sich der Nutzer aussuchen, ob er sich nur Informationen durchlesen, oder auf eine interaktive Weise Informationen über die Schule erfahren möchte. Hierbei werden zwei Buttons angezeigt, wobei beim Klick auf einem Button ein pdf-File heruntergeladen wird, welches Informationen enthält. Die Webseite enthält interaktive Elemente, gibt dem Nutzer die Möglichkeit die Medientechnik besser kennenzulernen und ist in verschiedenen Ebenen aufgebaut.
\item RE-M 2 Videos auf der Webseite\\
Die Videos sind auf der Webseite vorhanden.
Die gedrehten Videos sind auf der Webseite ersichtlich.
\item RE-M 3 Kontaktdaten\\
Die Kontaktdaten sind auf der Webseite vorhanden.\\
Bei Fragen können sich Nutzer, per E-Mail melden. Diese wird von schulinternen Personen betreut und ist jederzeit erreichbar. Außerdem befinden sich weitere Kontaktdaten wie Adresse und Telefonnummer auf der Webseite.
\item RE-M 4 Corporate Design\\
Ein Corporate Design ist erstellt.\\
Ein Corporate Design, welche die Farben, Schriftarten und das Logo der Webseite bestimmt, wird erstellt. Diese wird als Anhaltspunkt verwendet und ist als pdf-Datei verfügbar.
\item RE-M 5 Responsive Design\\
Die Webseite ist responsive ausgeführt.\\
Die Webseite ist responsive ausgeführt und kann auf mobilen und Desktop-Geräten bedient werden. 
\item RE-M 6 Drehbuch\\
Ein Drehbuch ist für das Tag der offenen Tür und Schülervideo erstellt.\\
Ein Drehbuch wird geschrieben. Dieses beinhaltet:\\
\begin{itemize}
\item Was gefilmt wird
\item Wie es gefilmt wird
\item Gesprochener Text 
\item Gebrauchtes Equipment
\item Dresscode
\end{itemize}
\item RE-M 7 Storyboard\\
Ein Storyboard für das Animationsvideo ist erstellt.\\
Ein Storyboard für das Animationsvideo wird erstellt um die Realisierung so einfach wie möglich zu gestalten.
\item RE-M 8 Genehmigung \\
Eine Genehmigung für das Filmen ist eingeholt.\\
Eine Drehgenehmigung wird beim Direktor eingeholt. Beim Filmen in den Klassenräumen wird ebenfalls um Erlaubnis gefragt, falls Schüler auf dem Material zu sehen sein sollten.
\item RE-M 9 Fragenkatalog\\
Ein Fragenkatalog ist erstellt.\\
Ein Fragenkatalog, welches die Fragen zum Interviewen beinhaltet, wird erstellt. Dabei ist zwischen den Fragen die an einen Medientechniklehrer und denen die an den Abteilungsvorstand gestellt werden, zu unterscheiden.
\item RE-M 10 Animationsvideo \\
Ein Animationsvideo ist erstellt.\\
Das Animationsvideo ist das erste Video, welches angezeigt wird. Dieses Video thematisiert allgemeine Fakten über die Schule und die Medientechnik. Das Video soll informativ, aber auch auflockernd wirken.
\item RE-M 11 Tag der offenen Tür\\
Ein Video ist am Tag der offenen Tür gedreht.\\
Am Tag der offenen Tür werden Eltern und Kinder nach ihrem ersten Eindruck befragt und gefilmt. Dabei wird auch einen Einblick in den Tag der offenen Tür geboten. Dieses Video beinhaltet außerdem zwei Interviews. Die Szenen werden dem Drehbuch entsprechend gefilmt.
\item RE-M 12 Interview\\
Zwei Interviews sind durchgeführt.\\
Ein Medientechniklehrer und der Abteilungsvorstand der Medientechnik werden interviewt. Die Fragen werden aus dem Fragenkatalog entnommen.
\item RE-M 13 Interviewer/in\\
Zwei Interviews sind durchgeführt.\\
Ein Medientechniklehrer und der Abteilungsvorstand der Medientechnik werden interviewt. Die Fragen werden aus dem Fragenkatalog entnommen.
\item RE-M 14 Voice-Over\\
Ein Voice-Over ist vertont.\\
Das Voice-Over beinhaltet den gesprochenen Text, welcher im Drehbuch festgelegt wurde. Dieser wird ins Video in den entsprechenden Stellen eingebunden.
\item RE-M 15 Schnitt \\
Das Video ist geschnitten.\\
Das Video wird dem Drehbuch entsprechend geschnitten.
\item RE-M 16 Color-Correction\\
Color-Correction ist durchgeführt.\\
Um die Belichtung und Farbe gekonnt in Szene zu setzen, wird Color-Correction durchgeführt. Hierbei wird darauf geachtet im natürlichen Bereich zu bleiben um das Video authentisch wie möglich wirken zu lassen.
\item RE-M 17 Interview mit Schülern\\
Ein Interview mit einem ehemaligen Schüler der Schule und einem Schüler der zweiten Klasse ist durchgeführt. \\
Ein ehemaliger Schüler und ein Schüler der zweiten Klasse müssen gemeinsam Fragen beantworten und die Antwort auf eine kleine Tafel schreiben und anschließend herzeigen und die Antwort begründen.
\item RE-M 18 FAQ\\
Ein FAQ ist erstellt.\\
Auf der Webseite befinden sich FAQ’s, um häufig gestellte Fragen zu beantworten.
\item RE-M 19 Visitenkarten\\
Visitenkarten sind erstellt und ausgedruckt.\\
Es werden Visitenkarten erstellt und am Tag der offenen Türe verteilt, um Aufmerksamkeit für die Diplomarbeit zu erwecken.
\item RE-M 20 Schulleitung \\
Die Schulleitung ist in die Diplomarbeit involviert.\\
Die Kommunikation, die Planung, das rechtliche Abklären und der Einfluss in das Content-Design, erfolgt mit der Schulleitung. 
\end{enumerate}

\section{Optionale Ziele}
Dabei wurden folgende optionale Ziele festgelegt:
\begin{enumerate}
\item RE-O 1 Interaktivität erhöhen \\
Mehrere interaktive Elemente sind vorhanden.\\
Die Webseite wird ausgebaut um mehrere interaktive Elemente einzubauen. Hierbei wird dem Nutzer die Komplexität der Medientechnik gezeigt.
\item RE-O 2 Quiz\\
Ein Quiz ist erstellt.\\
Auf der ersten Ebene wird der Nutzer aufgefordert ein Quiz zu machen. In diesem Quiz kann dieser herausfinden, ob er eher für die Medientechnik oder die Netzwerktechnik geeignet ist.
\end{enumerate}
\chapter{Design}
\renewcommand{\kapitelautor}{Autor: Niklas Kienreich}

\section{Hintergedanken}
Genau, wie bei allen anderen Aspekten, dieser Diplomarbeit, war der Hintergedanke des Designs, es einladend, fröhlich und angemessen der Zielgruppe zu gestalten.

\section{Richtlinien}
\section{Corporate Design}
\subsection{Farben}
Bei den Farben wurde zum einen, ein helles Grün, zum anderen ein knalliges Orange gewählt. Als Akzentfarbe wurde ein helles Grau in das Corporate Design aufgenommen.

\subsubsection{Grün}
Wenn etwas ergrünt, ergibt sich Hoffnung auf neues Leben. Diese Farbe ist sehr naturbelassen und die Farbpsychologie lehrt, dass Grün Lebensfreude und Jugend vermittelt. Grün kann zwar auch Gift und somit eine Warnung signalisieren, durch den Kontext der Webseite ist dies aber unwahrscheinlich. \footnote{\label{foot:1} http://www.grafixerin.com/bilder/Farbpsychologie.pdf [Zugriff: 17.03.2018]}

\subsubsection{Orange}
Bei Orange handelt es sich um eine sehr auffällige Farbe. Sie ist eine warme, energiegeladene Farbe und kann aufregend, bis sogar spaßig auf einen Menschen wirken. Durch ihre Auffälligkeit kann sie aufdringlich oder extrovertiert wirken. Orange kann auch billig wirken, da der Mensch kaum etwas Hochwertiges kennt, dass diese Farbe besitzt. \footnote{\label{foot:1} http://www.grafixerin.com/bilder/Farbpsychologie.pdf [Zugriff: 17.03.2018]}

\subsubsection{Grau}
\subsection{Typografie}
\subsubsection{Parisien Night Oblique}
Diese Schrift wurde wegen ihrer dynamischen Linienführung gewählt, welche an das Logo erinnert.
\\
Verwendung findet die Schrift in den Überschriften der Website.

\subsubsection{Calibri}
Calibri ist eine serifenlose Schrift. Sie ist dadurch auf Bildschirmen mit geringerer Auflösung leicht zu lesen. Sie wirkt durch ihre abgerundeten Ecken freundlicher, als manch andere Schriftart. \footnote{\label{foot:2} https://www.typografie.info/3/Schriften/fonts.html/calibri-r61/ [Zugriff: 18.03.2018]}
\\
Calibri wird in den Texten der Website und in den Textslides der Videos verwendet.

\subsection{Adobe Illustrator}
\subsubsection{Vectorgrafiken}

\subsection{Logo}
\subsubsection{Logo-Arten}
Bei Logos unterscheidet man zwischen vier verschiedenen Arten. \footnote{\label{foot:3} http://www.webmasterpro.de/design/article/logodesign-arten-von-logos.html [Zugriff: 18.03.2018]} \footnote{\label{foot:4} http://www.sanseg.de/logodesign/zeichenmarken/ [Zugriff: 18.03.2018]}
\\
\\
Zum einen gibt es sogenannte Bildmarken. Hierbei handelt es sich um eine einfache Abbildung oder Illustration. Empfehlenswert ist solch ein Logo jedoch eher für größere Firmen und Unternehmen, da der Betrachter beim Erblicken des Logos, die Verbindung mit jenem Unternehmen herstellen können soll. Dies ist bei unbekannten Marken unwahrscheinlich und somit kontraproduktiv. Wenn ein Unternehmen auf eine Bildmarke umsteigt, gibt es oft eine Einführungsphase, in der das neue Logo anfangs mit Beschriftung gezeigt wird.
\\
\\
Eine weitere Logo-Art ist die Wortmarke. Sie zeigt einfach den Namen, des Unternehmens. Aber auch hier muss man sich Gedanken machen, wie das Logo auf den Betrachter wirkt. Neben der Farbe des Schriftzugs, sind typografische Aspekte auch wichtig. Darauf wird später in der Umsetzung des Logos genauer eingegangen.
\\
\\
Zeichenmarken sind den Wortmarken sehr ähnlich, denn sie beruhen auch auf der Typografie. Der Unterschied liegt darin, dass es sich hier nicht um Worte im Logo handelt, sondern um Abkürzungen, einzelne Buchstaben oder Zahlen.
\\
\\
Bei Wort-Bild-Marken, oder auch Kombinationen, handelt es sich, wie der Name schon verrät, um eine Kombination aus Bildmarke und entweder Wortmarke, oder Zeichenmarke. Bei dieser Art des Logos handelt es sich um die weit verbreitetste Variante. Sie ist auch am besten für kleine und mittelgroße Unternehmen geeignet.
\\
\\
Auch Keyvisuals können genutzt werden. Sie sind nicht an das Logo gebunden und können somit davon abweichen und eigene Charakteristiken veranschaulichen. Ein Keyvisual ist dann gut, wenn auf ihm der Name des Unternehmens nicht vertreten ist, der Betrachter es aber trotzdem ohne viel nachzudenken zuordnen kann.

\subsubsection{Zweck}
\subsubsection{Symbolik}
\subsubsection{Umsetzung}
\subsection{Visitenkarten}
\subsection{Designelemente}

\section{Webseite}
\subsection{Nutzen}
\subsection{Bildquellen}
\subsection{Frameworks und Libraries}
\subsubsection{Foundation}
\subsubsection{CreateJS}
\subsubsection{Timeline.js}
\subsection{Responsiveness}
\subsection{Storytelling}
\subsubsection{Story}
\subsubsection{Figur}

\subsection{Inhalt}
\subsubsection{Startseite}

\subsubsection{Interaktive Elemente}



\chapter{Video}
\section{Allgemeines}
\renewcommand{\kapitelautor}{Autor: Kerstin Schön}
Dieses Kapitel befasst sich mit der Planung, der Durchführung und der Bearbeitung der Videos. 

\subsection{Storyboard}
\renewcommand{\kapitelautor}{Autor: Kerstin Schön}
Um die korrekte Ausführung der Videos zu gewährleisten, wurden sogenannte Storyboards angefertigt. Diese wurden mit der Online Plattform "StoryboardThat" erstellt und schließlich als PDF - Format exportiert.
Im Storyboard wurden die Szenen bildlich dargestellt. Dies unterstützte das Team insofern im weiteren Verlauf, um grobe Fehler zu vermeiden.
Das Online - Tool ermöglichte es zwischen verschiedensten Szenen, Charaktären und Kategorien auszuwählen, wodurch einfach, verschiedenste Szenen erstellt werden konnten.
 
\subsection{Fragenkatalog}
\renewcommand{\kapitelautor}{Autor: Kerstin Schön}
Zusätzlich zum Storyboard wurden für das jeweilige Video die entsprechenden Fragen erstellt. Hierbei war es wichtig die Fragen an die Zielgruppe anzupassen. 

\subsection{Setting}
Beim Aufbau des Sets wurde zunächst auf das Storyboard referenziert, welche Situation man vorfindet. 

\subsection{Kamera}
%https://www.canon.at/for_home/product_finder/cameras/digital_slr/technologies_features/cmos_sensor.aspx
%https://www.itwissen.info/CMOS-Sensor-CMOS-sensor.html
Zum Drehen der Videos kam die Spiegelreflexkamera Canon EOS 70D zum Einsatz. Die Kamera Canon EOS 70D hat den CMOS Sensor verbaut, welcher im Gegensatz zu dem CCD (charge-coupled device) - Sensor einige Vorteile mit sich bringt. Vereinfacht gesagt, sind CMOS bzw. CCD - Sensoren "lichtempfindliche Bauteile, die das auf sie fallende Licht in Spannung umwandeln".  

\subsection{Licht}
Um den Schauplatz genügend auszuleuchten, verwendeten wir für alle vier Videos LED Scheinwerfer. -> softbox 

\subsection{Mikrofone}
%https://shuredeutschland.wordpress.com/2013/02/08/mikrofongrundlagen-richtcharakteristik/
%http://www.shure.at/supportdownload/tipps_grundlagen/mikrofone/mikrofone-richtcharakteristik
%https://www.delamar.de/mikrofon/richtcharakteristik-mikrofon-22647/
Bevor man sich für ein Mikrofon entscheidet, müssen einige Dinge beachtet werden, um das optimalste Ergebnis zu erreichen. Es gibt Mikrofone mit einer Kugelcharakteristik, die sich z.B. für Aufnahmen gut eignen, wenn der Schall von allen Seiten aufgenommen werden soll. Will man möglichst gerine Hintergrundgeräusche erzielen, eignet sich optimal ein Mikrofon mit einer Nierencharakteristik. \\
Bei der Richtcharakteristik eines Mikrofons kommt es grundsätzlich darauf an, aus welcher Richtung man den Schall aufnehmen möchte. Eine Kugelcharakteristik eignet sich dann z.B. gut, wenn man mehrere Leute


verwendung von zoom bzw. einmal mit handy vorteile und nachteile 

\section{Interview mit dem Abteilungsvorstand}
\renewcommand{\kapitelautor}{Autor: Kerstin Schön}
\subsection{Idee}
\renewcommand{\kapitelautor}{Autor: Kerstin Schön}
\section{Tag der offenen Tür}
\renewcommand{\kapitelautor}{Autor: Kerstin Schön}
\subsection{Idee}
\renewcommand{\kapitelautor}{Autor: Kerstin Schön}
\section{Video mit einem Absolvent}
\renewcommand{\kapitelautor}{Autor: Kerstin Schön}
\subsection{Idee}
\renewcommand{\kapitelautor}{Autor: Kerstin Schön}

\section{Schnitt}
\renewcommand{\kapitelautor}{Autor: Kerstin Schön}

\section{Herausforderungen}
\renewcommand{\kapitelautor}{Autor: Kerstin Schön}
\subsection{Tiefenschärfe}
\renewcommand{\kapitelautor}{Autor: Kerstin Schön}
\subsection{Farbkorrektur}
\renewcommand{\kapitelautor}{Autor: Kerstin Schön}
\subsection{Audio}
\renewcommand{\kapitelautor}{Autor: Kerstin Schön}
\subsection{Vertonung}
\renewcommand{\kapitelautor}{Autor: Kerstin Schön}

\section{Videoexport}
\renewcommand{\kapitelautor}{Autor: Kerstin Schön}
\subsection{Browserkompatibilität}
\renewcommand{\kapitelautor}{Autor: Kerstin Schön}


\section{Interview mit einem Fachmann}
\subsection{Idee}
Wir wollten ein Interviewvideo mit einem passionierten Medientechniklehrer drehen, um den Usern unserer Plattform, einen tiefen Einblick in die Medientechnik zu bieten.

\subsection{Ziel}

\subsection{Dreh}

\subsection{Schnitt}

\subsection{Green Screen}

\subsection{Farbkorrektur}



\chapter{Animationsvideo}

\section{Idee}
\renewcommand{\kapitelautor}{Autor: Niklas Kienreich}
Das Animationsvideo sollte als Eyecatcher für die Zielgruppe dienen. Der Gedanke war, Allgemeines über die HTL auf möglichst witzige, ansprechende Weise zu vermitteln. Die Animation schafft das besser, als die restlichen Videos, da man sich bei den Interviews auf eine lockere, aber doch ernste, zielführende Gesprächsführung verlassen hat. Eine Frage die man sich nun aber stellen musste war, wie man das Video animiert. Neben vier Interview Videos, die nicht nur gedreht, sondern auch geschnitten und farbkorrigiert werden mussten und der Website, war für so ein kleines Team von drei Personen nicht allzu viel Zeit eingeplant. Nach überraschend kurzer Recherche, war nicht nur eine einfache, sondern auch ansprechende Lösung gefunden. In den letzten Jahren haben immer mehr Kanäle auf Youtube bei unserer Zielgruppe großes Interesse geweckt. Laut Youtube exestiert einer der beliebteren Kanäle seit dem 30.08.2014 und hat seitdem fast 6 Millionen Abonnenten und 866 Millionen Aufrufe gesammelt. 
Ihre Videos sind kurze Animationen im vereinfachten Stil. Also keine flüssigen Bewegungen, sondern mehr sprunghafte Frames mit einem Voice-over. Es erinnert an eine digitale Version eines so genannten „Draw my life“.
\section{Ziel}
\renewcommand{\kapitelautor}{Autor: Niklas Kienreich}
Das Ziel war es, mit dem zwar kürzesten, aber unterhaltsamsten Video, die wichtigsten, grundlegenden Informationen zu Höheren Technische Bundeslehranstalten zu bieten. Die detaillierten Informationen zu der Schule selbst, sind im Interviewvideo mit dem Abteilungsvorstand zu finden. Die User sollten nach den beiden Videos zuerst wissen, ob sie an solch einer Schule überhaupt Interesse haben. Ein Wissen das essenziell ist, bevor man sich auf die Frage stürzt, ob die Medientechnik, oder gar irgendein anderer technischer Fachbereich für einen geeignet ist.

\section{Drehbuch}
\renewcommand{\kapitelautor}{Autor: Niklas Kienreich}

\subsection{Brainstorming}
Bevor das Drehbuch geschrieben werden konnte, musste ein Brainstorming durchgeführt werden. Hierbei wurden zuerst die Drehbücher beziehungsweise Fragen, der anderen Videos herangezogen, um zu ermitteln, welche Fragen noch nicht geklärt wurden, oder welche Fragestellung man nochmal genau herausheben sollte.
\leavevmode \\
Das Brainstorming brachte folgende Ergebnisse:
\begin{itemize}
\item Was genau ist eine HTL?
\item Welche Fachbereiche gibt es und wie teilen sie sich auf?
\item Was lernt man an der Schule unter anderem?
\item Was für Möglichkeiten bieten sich einem nach der Schule?
\end{itemize}

\subsection{Wirkung}
Da das Animationsvideo als einleitendes Video genutzt wird, ist es sehr wichtig, wie es auf den Benutzer wirkt. Um das Video freundlich, einladend und möglichst natürlich wirken zu lassen, wurde beschlossen, dass es in dem Video um einen Schüler geht, mit dem sich der User identifizieren können soll.

\section{Audioaufnahme}

\subsection{Mikrofon}

\subsection{Software}

\section{Schnitt}

\section{Umsetzung}


\appendix

\chapter{Anhang 1\label{chap:Anhang-1}}

was auch immer: technische Dokumentationen etc.

Zusätzlich sollte es geben: 
\begin{itemize}
\item Abkürzungsverzeichnis
\item Quellenverzeichnis (hier: Bibtex im Stil plaindin)
\end{itemize}
\printindex{}

%% Flattersatz -- damit werden die langen URLs besser umgebrochen
\raggedright %% eventuell auskommentieren
\bibliographystyle{plaindin}%Alternative unsrtdin - Nummern im Text aufsteigend
\bibliography{diplom}


\cleardoublepage
\newcommand{\Messbox}[2]{%Parameters: #1=Breite, #2=Hoehe
\setlength{\unitlength}{1.0mm}%
\begin{picture}(#1,#2)%
\linethickness{0.05mm}%
\put(0,0){\dashbox{0.2}(#1,#2)%
{\parbox{#1mm}{%
\centering\footnotesize  
%{\bf MESSBOX}\\
Breite $ = #1 {\textrm\ mm}$\\
Höhe $ = #2 {\textrm\ mm}$
}}}\end{picture}
}
\begin{center} {\Large --- Druckgröße kontrollieren! ---}
\bigskip

\Messbox{100}{50} % Angabe der Breite/Hoehe in mm
\bigskip

{\Large --- Diese Seite nach dem Druck entfernen! ---}
\todo{Diese Seite nach dem Druck entfernen!}
\end{center} 
\end{document}
